% ============================================================
% INTRODUÇÃO
% ============================================================

% -------------------------------------------------------
\begin{frame}
    \vspace{.8cm}
    \titlepage
\end{frame}

% -------------------------------------------------------
\section{Introdução}
% -------------------------------------------------------

\begin{frame}{Motivação}
    \begin{itemize}
        \item A literatura sobre o impacto de estradas na economia tem se concentrado no \textbf{movimento de bens}.
        \item Pouca atenção foi dada ao papel das estradas como \textbf{facilitadoras da migração} interna (movimento de pessoas).
        \item Este artigo pergunta:
        \begin{enumerate}
            \item Estradas \textbf{promovem migração}?
            \item Qual parcela dos ganhos de bem-estar das estradas decorre de \textbf{integração de mercados de bens} vs.\ \textbf{integração de mercados de trabalho}?
        \end{enumerate}
    \end{itemize}
\end{frame}

\begin{frame}{Contribuições}
    \begin{block}{Contribuição 1 -- Evidência empírica}
        Estradas promovem migração. Uso de variação \textbf{plausivamente exógena} na localização da malha rodoviária (construção de Brasília).
    \end{block}

    \vspace{.4cm}

    \begin{block}{Contribuição 2 -- Decomposição do bem-estar}
        Decompõe os ganhos de bem-estar das rodovias entre os canais de \textbf{comércio} (76\%) e de \textbf{migração} (24\%). Migração custosa é responsável por heterogeneidade espacial significativa.
    \end{block}
\end{frame}

\begin{frame}{Sumário executivo dos resultados}
    \begin{itemize}
        \item A melhoria das estradas aumentou o bem-estar em \textbf{2,8\%}.
        \item \textbf{76\%} desse ganho decorreu da redução de custos de comércio; \textbf{24\%}, de custos de migração.
        \item Na ausência de migração custosa, os ganhos seriam uniformes no espaço. Com ela, a amplitude regional fica entre \textbf{1\% e 15\%}.
        \item Exercício contrafactual: a rede construída em torno de Brasília gerou bem-estar \textbf{0,3\% maior} que uma rede alternativa centrada no Rio de Janeiro.
    \end{itemize}
\end{frame}
